\section{TODD：早期棄却法の改良}

前回実装した早期棄却法を，それまで最速であった基底再利用法と組み合わせた手法を実装した．
目的は，完全な $\chi$ 行列の更新や行空間基底の再構築を行う前に，
削減不可能な候補を低コストに棄却し，基底更新や零空間計算に進む候補数を減らすことである．
本手法は，所属判定，早期棄却，基底再利用の 3 つを組み合わせたものである．

\subsection{所属判定}

元の TODD 実装では，各候補ペア $(c_1,c_2)$ に対して
$\chi$ 行列 $M$ の零空間基底を求め，
\[
    y[c_1] \oplus y[c_2] = 1
\]
を満たすベクトル $y \in \operatorname{Null}(M)$ が存在するかを調べる．

一方，基底再利用法では，零空間計算の前に行空間への所属判定を行う．
候補ペア $(c_1,c_2)$ に対して，
\[
    v = e_{c_1}+e_{c_2}
\]
とおく．
ここで，$v$ は $c_1$ 番目と $c_2$ 番目のみが $1$ であるベクトルである．
任意のベクトル $y$ に対して，
\[
    v \cdot y = y[c_1] \oplus y[c_2]
\]
が成り立つ．

行列 $M$ の行空間と零空間は直交するため，
\[
    v \in \operatorname{Row}(M)
\]
であれば，任意の $y \in \operatorname{Null}(M)$ に対して
\[
    v \cdot y = 0
\]
となる．
したがって，
\[
    e_{c_1}+e_{c_2} \in \operatorname{Row}(M)
\]
であれば，
\[
    y[c_1] \oplus y[c_2] = 1
\]
を満たす零空間ベクトルは存在せず，候補ペア $(c_1,c_2)$ は削減不可能であると判定できる．

本実装において，行空間基底は掃き出し済みの形で保持されている．
したがって，零空間基底を構成すること自体は可能であるが，
その場合には自由変数に対応する零空間ベクトルを生成し，
各ベクトルに対して $y[c_1]\oplus y[c_2]$ を確認する必要がある．
一方，所属判定では，保持している行空間基底をそのまま用いて，
候補ベクトル $e_{c_1}+e_{c_2}$ の 1 本だけを簡約すればよい．
行空間基底の本数を $r$，列数を $m$ とすると，所属判定の計算量は
\[
    O(rm)
\]
このため，漸近的な計算量の差だけではなく，
零空間ベクトルの生成やメモリ確保を避けられる点で，実装上は定数倍で軽くなる可能性がある．

\subsection{早期棄却}

早期棄却では，完全な $\chi$ 行列 $M$ に対する判定を行う前に，
その一部の行からなる小型 $\chi$ 行列 $M'$ を用いる．
$M'$ は $M$ の行部分集合であるため，
\[
    \operatorname{Row}(M') \subseteq \operatorname{Row}(M)
\]
が成り立つ．
したがって，
\[
    e_{c_1}+e_{c_2} \in \operatorname{Row}(M')
\]
であれば，
\[
    e_{c_1}+e_{c_2} \in \operatorname{Row}(M)
\]
も成り立つ．
この場合，完全な $\chi$ 行列を用いなくても候補ペアが削減不可能であることを安全に判定できる．

本手法は，完全な $\chi$ 行列を更新する前に，
ランダムに選んだ $p\%$ の行から小型 $\chi$ 行列を構成し，所属判定を行う．
小型行列で削減不可能と判定できた候補については，
完全な $\chi$ 行列の更新，行空間基底の更新，および零空間計算を行わずに棄却する．

\subsection{基底再利用}

基底再利用法では，候補ごとに完全な $\chi$ 行列の行空間基底をゼロから構築するのではなく，
前回の候補で得られた $\chi$ 行列の状態と行空間基底を保持する．
次の候補では，前回の $\chi$ 行列から変化した行のみを調べる．

変化した行が現在の行空間基底を構成している基底行を含まない場合，
既存の基底を再利用し，新たに有効になった行だけを追加する．
一方，変化した行が基底行を含む場合には，
既存の基底が現在の行空間を正しく表さない可能性があるため，行空間基底を再構築する．

今回の手法では，この基底再利用の前段に早期棄却を挿入する．
これにより，完全な $\chi$ 行列の状態更新自体を行わない候補を増やし，
基底更新や再構築に進む候補数を削減する．

\subsection{手法の流れ}

各候補ペア $(c_1,c_2)$ に対する処理は以下の通りである．

\begin{enumerate}
    \item 候補ペア $(c_1,c_2)$ に対して，
    \[
        x = A[:,c_1] + A[:,c_2]
    \]
    を計算する．

    \item 完全な $\chi$ 行列の状態および行空間基底はまだ更新しない．

    \item 候補に対応する $\chi$ 行列の一部の行をランダムに選び，小型 $\chi$ 行列 $M'$ を構成する．

    \item 小型 $\chi$ 行列 $M'$ に対して，
    \[
        e_{c_1}+e_{c_2} \in \operatorname{Row}(M')
    \]
    であるかを判定する．

    \item 所属する場合，候補ペアは削減不可能であるため，完全な $\chi$ 行列の更新を行わずに棄却する．

    \item 所属しない場合のみ，完全な $\chi$ 行列の状態と行空間基底を更新する．
    変化した行が基底行を含まなければ既存の基底を再利用し，含む場合は基底を再構築する．

    \item 完全な $\chi$ 行列に対して再度所属判定を行い，
    \[
        e_{c_1}+e_{c_2} \in \operatorname{Row}(M)
    \]
    であれば候補ペアを棄却する．

    \item 完全な行空間所属判定でも棄却できない場合のみ，
    零空間計算を行い，削減に用いるベクトル $y$ を探索する．
\end{enumerate}

